\documentclass[12pt,spanish,letterpaper,oneside]{article}
\def\documenttitle {Utilidad actual de las generaciones de derechos humanos}
\def\documentsubtitle {Generaciones: protección y garantía}
\def\documentsubject {Generaciones de los derechos humanos: protección y garantía}
\def\documentauthor {Martin Jonathan de la Cruz}
\def\coursename {Antecedentes de los derechos humanos}
\def\coursecode {005}
\def\universityname {Universidad Abierta y a Distancia de México}
\def\universityfaculty {Licenciatura en Derecho}
\def\universitydepartment {Antecedentes de los derechos humanos}
\def\universitydepartmentimage {departamentos/UnADM}
\def\universitydepartmentimagecfg {height=1.57cm}
\def\universitylocation {Ciudad de México}
\def\authortable {\begin{tabular}{ll}
  \textbf{Alumno:} & Martin Jonathan de la Cruz \\
  Matricula: & ES2611202040 \\
  Figura docente: & Daniel Alexis Lozano Ortega \\
  Semestre/Bloque: & 2 / 2 \\
  Semana: & 3 \\
  Fecha: & 18 de octubre de 2026
\end{tabular}}
\usepackage{iftex}
\ifLuaTeX
  \let\pdfpagewidth\pagewidth
  \let\pdfpageheight\pageheight
\fi
\input{template}
\usepackage{fontspec}
\setmainfont{arial.ttf}[Path=C:/Windows/Fonts/,BoldFont=arialbd.ttf,ItalicFont=ariali.ttf,BoldItalicFont=arialbi.ttf]
\setsansfont{arial.ttf}[Path=C:/Windows/Fonts/,BoldFont=arialbd.ttf,ItalicFont=ariali.ttf,BoldItalicFont=arialbi.ttf]
\renewcommand{\familydefault}{\sfdefault}
\setcitestyle{authoryear,open={(},close={)}}
\usepackage[most]{tcolorbox}
\usepackage{fvextra}
\usepackage{attachfile2}
\usepackage{tikz}
\usepackage{eso-pic}
\definecolor{unadmGreen}{HTML}{5F8F3A}
\definecolor{unadmGreenDark}{HTML}{174A3A}
\definecolor{unadmPaper}{HTML}{F6F7F2}
\attachfilesetup{color={0.373 0.561 0.227},print=false}
\newcommand{\foroCopyButton}[1]{%
  \textattachfile[appearance={},description={Participacion en texto},mimetype=text/plain]{#1}{%
    \colorbox{unadmGreen}{\small\textcolor{white}{\,Abrir texto\,}}}}
\newcommand{\insertcoverwatermark}{%
  \AddToShipoutPictureBG*{%
    \begin{tikzpicture}[remember picture,overlay]
      \node[opacity=0.16,inner sep=0pt] at (current page.center)
        {\includegraphics[width=0.72\paperwidth]{img/departamentos/UnADM.pdf}};
    \end{tikzpicture}%
  }%
}
\newtcolorbox{forobox}[1][]{enhanced,breakable,colback=unadmPaper,
  colframe=unadmGreen,boxrule=0pt,leftrule=3pt,arc=1.2mm,
  left=4mm,right=4mm,top=3mm,bottom=3mm,
  fonttitle=\bfseries\color{white},coltitle=white,colbacktitle=unadmGreenDark,
  attach boxed title to top left={xshift=4mm,yshift=-2mm},
  boxed title style={colframe=unadmGreenDark,arc=0.6mm},#1}
% Foro (#29): cada archivo es el texto integro seleccionable de una intervencion.
\begin{document}
\insertcoverwatermark
\templatePortrait
\templatePagecfg
\templateFinalcfg
\newgeometry{margin=2.5cm}
\fontsize{12}{14.5}\selectfont
\onehalfspacing

\section{Introducción}
La clasificación de los derechos humanos en generaciones permite reconocer el contexto en que distintas libertades y exigencias sociales adquirieron protección jurídica. Su utilidad actual depende de una condición: comprender esos antecedentes sin convertirlos en una jerarquía que separe derechos interdependientes. En México, el contraste entre historia y garantías obliga a preguntar qué autoridades deben actuar y mediante qué medidas pueden hacerse efectivos los derechos.

\section{Historia, interdependencia y exigibilidad}
Las generaciones relacionan demandas de libertad, igualdad y solidaridad con contextos históricos distintos; no implican que los derechos posteriores sustituyan a los anteriores \citep[pp.~72--76]{reyCantor2022Foro}. Su explicación histórica debe distinguirse de las obligaciones actuales: el artículo 1 constitucional exige promover, respetar, proteger y garantizar los derechos con interdependencia e indivisibilidad \citep[art.~1]{cpeumForoS3}. Identificar responsables y medidas de protección permite pasar de esa clasificación a la exigibilidad.

\subsection{Protección integral de los derechos}
\begin{forobox}[title={Generaciones y protección integral}]
\hfill\foroCopyButton{UnADM/licenciatura-en-derecho-unadm/antecedentes-de-los-derechos-humanos-lde/foro-participacion-Actividad-6.txt}\par\medskip
\VerbatimInput[breaklines=true,breakanywhere=true,fontsize=\footnotesize,baselinestretch=1,breakindent=1.27cm,breaksymbolleft={},breaksymbolright={}]{foro-participacion-Actividad-6.txt}
\end{forobox}

La interdependencia de los derechos también atraviesa la desigualdad y el daño ambiental: ambas situaciones dificultan que las libertades formales se ejerzan en condiciones reales. Distinguir el origen histórico de los derechos de las obligaciones presentes evita reducir su protección a una cronología.

\subsection{Protección conjunta y cumplimiento efectivo}
La posición a la que se responde subraya que ninguna generación es superior a otra y que el deterioro ambiental exige protegerlas conjuntamente. Precisar cómo se mide el cumplimiento estatal permite llevar ese argumento a sus consecuencias prácticas.
\begin{forobox}[title={Interdependencia y obligaciones estatales}]
\hfill\foroCopyButton{UnADM/licenciatura-en-derecho-unadm/antecedentes-de-los-derechos-humanos-lde/foro-replica-Actividad-6.txt}\par\medskip
\VerbatimInput[breaklines=true,breakanywhere=true,fontsize=\footnotesize,baselinestretch=1,breakindent=1.27cm,breaksymbolleft={},breaksymbolright={}]{foro-replica-Actividad-6.txt}
\end{forobox}

La respuesta distingue abstenciones y medidas positivas: evitar una interferencia no siempre basta para asegurar el ejercicio de un derecho. La pregunta sobre indicadores de cumplimiento permite contrastar el argumento con acciones verificables, sin presentar la clasificación histórica como prueba de protección efectiva \citep[art.~1]{cpeumForoS3}.

\clearpage
\section{Conclusiones}
Las generaciones explican la aparición histórica de distintas demandas de libertad, igualdad y solidaridad, pero no determinan la importancia relativa de los derechos. A mi juicio, su utilidad jurídica radica en situar esas demandas sin perder de vista que su cumplimiento depende de obligaciones estatales concurrentes. La consecuencia es examinar, en cada caso, tanto el derecho afectado como las acciones y garantías necesarias para ejercerlo: la mera clasificación cronológica no demuestra protección efectiva.\footnote{Uso de IA: GitHub Copilot para organización conceptual y revisión lingüística; la valoración jurídica se contrasta con las fuentes citadas.}

\clearpage
\begin{thebibliography}{3}
\bibitem[Camara de Diputados(2026)]{cpeumForoS3}
Cámara de Diputados. (2026). \emph{Constitución Política de los Estados Unidos Mexicanos} (art.~1). \url{https://www.diputados.gob.mx/LeyesBiblio/pdf/CPEUM.pdf}
\bibitem[Rey Cantor et~al.(2022)]{reyCantor2022Foro}
Rey Cantor, E., Rey Anaya, G. A., \& Rey Anaya, A. M. (2022). \emph{Las generaciones de los derechos humanos: Libertad, igualdad y fraternidad} (10.\textsuperscript{a} ed.). Universidad Libre.
\bibitem[UK Parliament(s.~f.)]{parlamentoCarta1215Foro}
UK Parliament. (s.~f.). \emph{Magna Carta}. \url{https://www.parliament.uk/about/living-heritage/evolutionofparliament/originsofparliament/birthofparliament/overview/magnacarta/}
\end{thebibliography}
\end{document}